\documentclass[10pt,letterpaper]{article}

\usepackage[T1]{fontenc}
\usepackage[utf8]{inputenc}
\usepackage{newtxtext,newtxmath}
\usepackage[letterpaper,margin=.65in,top=.7in,bottom=.65in]{geometry}
\usepackage{microtype}
\usepackage{xcolor}
\usepackage{tabularx}
\usepackage{booktabs}
\usepackage{array}
\usepackage{enumitem}
\usepackage[most]{tcolorbox}
\usepackage{titlesec}
\usepackage{fancyhdr}
\usepackage{hyperref}

\definecolor{cvblue}{HTML}{243746}
\definecolor{cvred}{HTML}{7A2E2E}
\definecolor{cvsand}{HTML}{F0EBE1}
\definecolor{cvink}{HTML}{202020}

\hypersetup{
  colorlinks=true,
  linkcolor=cvred,
  urlcolor=cvblue,
  pdftitle={Cromwell's Victory: Rupert Strikes at Noon},
  pdfsubject={Experimental alternate scenario for the Battle of Marston Moor},
  pdfauthor={Daniel J. Berger}
}

\titleformat{\section}
  {\sffamily\bfseries\Large\color{cvblue}\raggedright}
  {}{0pt}{}[\color{cvred}\titlerule]
\titlespacing*{\section}{0pt}{.8em}{.45em}
\titleformat{\subsection}
  {\sffamily\bfseries\normalsize\color{cvred}\raggedright}
  {}{0pt}{}
\titlespacing*{\subsection}{0pt}{.55em}{.2em}

\newtcolorbox{scenarioBox}[1]{
  enhanced,breakable,
  colback=cvsand!48,colframe=cvblue!72,
  boxrule=.55pt,arc=1pt,
  left=7pt,right=7pt,top=4pt,bottom=4pt,
  before skip=5pt,after skip=6pt,
  title={\sffamily\bfseries\small #1},
  coltitle=white,colbacktitle=cvblue
}

\newtcolorbox{redBox}[1]{
  enhanced,breakable,
  colback=cvsand!35,colframe=cvred!82,
  boxrule=.55pt,arc=1pt,
  left=7pt,right=7pt,top=4pt,bottom=4pt,
  before skip=5pt,after skip=6pt,
  title={\sffamily\bfseries\small #1},
  coltitle=white,colbacktitle=cvred
}

\newcolumntype{Y}{>{\raggedright\arraybackslash}X}
\newcommand{\turn}[1]{{\sffamily\bfseries\color{cvred}Turn #1}}

\setlength{\parindent}{0pt}
\setlength{\parskip}{.38em plus .1em minus .06em}
\renewcommand{\arraystretch}{1.13}
\setlist[itemize]{leftmargin=1.2em,itemsep=.22em,topsep=.2em,parsep=0pt}
\setlist[enumerate]{leftmargin=1.45em,itemsep=.28em,topsep=.2em,parsep=0pt}

\pagestyle{fancy}
\fancyhf{}
\fancyhead[L]{\sffamily\small\color{cvblue}CROMWELL'S VICTORY}
\fancyhead[R]{\sffamily\small\color{cvblue}RUPERT STRIKES AT NOON}
\fancyfoot[C]{\sffamily\small\thepage}
\renewcommand{\headrulewidth}{.35pt}
\setlength{\headheight}{14pt}

\begin{document}

\begin{center}
  {\color{cvred}\rule{\textwidth}{1.5pt}}\par\vspace{8pt}
  {\sffamily\bfseries\fontsize{29}{32}\selectfont\color{cvblue}
    RUPERT STRIKES AT NOON}\\[2pt]
  {\sffamily\large\color{cvred}An alternate scenario for \textsc{Cromwell's Victory}}\\[4pt]
  {\sffamily\small Daniel J. Berger \textbullet{} \today}
\end{center}

\begin{redBox}{Experimental Scenario}
This scenario has not yet been balanced by repeated play. It is meant to test
the battle Prince Rupert wanted: Newcastle's York infantry reaches the moor in
time for a noon attack, before the Allied foot and artillery have returned and
formed their line.
\end{redBox}

\section{Historical Situation}

Early on 2 July 1644 the Allied army began marching toward Tadcaster. Its horse
remained behind as a rearguard, but much of its infantry, artillery, and baggage
was already on the road when Rupert appeared in strength. The Allied commanders
hurriedly recalled the column and began rebuilding their line on Marston Hill.

Newcastle joined Rupert during the morning, but Eythin did not bring the York
infantry onto the moor until late afternoon. By then the opportunity had passed:
the Allied deployment was substantially complete between 2 and 3 p.m., while
Newcastle's infantry arrived between 4 and 5 p.m. The historical evidence and
timing are summarized in Historic England's
\href{https://historicengland.org.uk/content/docs/listing/battlefields/marston-moor/}{battlefield report}
and the \href{https://bcw-project.org.uk/military/english-civil-war/northern-england/battle-of-marston-moor}{BCW Project account}.

Here Eythin gets the York regiments moving promptly. Rupert has his full army
at noon and attacks while the Allied horse is still covering an infantry column
that is arriving piecemeal.

\begin{scenarioBox}{What the Scenario Is Testing}
The Royalists receive a brief superiority in formed foot and artillery, but the
advantage erodes every half-hour. The Royalist player must press the Allied
rearguard or seize the forming ground before Leven's large infantry contingent
arrives. The Allied player must trade space for time without allowing the horse
to be trapped or the assembly area to be overrun.
\end{scenarioBox}

\subsection{Length}

The scenario lasts 12 half-hour turns, beginning at noon. Advance the time by
30 minutes each turn; the scenario ends at 6 p.m. after Turn 12.

\section{Set Up}

Set up the map, Game Turn Marker, and train hexes normally. All counters begin
on their ordered sides unless stated otherwise.

\subsection{Royalist Army}

Set up \textbf{every counter of Rupert's and Newcastle's forces} in the set-up
hex printed on its counter. This includes Newcastle, Eythin, all seven Newcastle
Foot, all Northern Horse, and all four Royalist artillery pieces. The earlier
arrival has already happened; the Royalist army is assembled when play begins.

\subsection{Allied Rearguard}

Set up the following counters in their printed set-up hexes:

\begin{itemize}
  \item every Allied Heavy Horse, Light Horse, and Dragoon unit;
  \item Manchester, Cromwell, Leven, Fairfax, and Thomas Fairfax; and
  \item no Allied Foot or artillery.
\end{itemize}

Place all Allied Foot and all six Allied artillery pieces off-map, sorted into
the groups shown below.

\section{Allied Arrival Schedule}

Arriving pieces enter during their scheduled Allied March Phase through the
\textbf{southwest entry area}. This consists of south-edge hexes
0101--0108 and west-edge hexes 0101--0801.

The first map hex entered costs its normal terrain cost; there is no additional
cost for crossing the map edge. A piece need not enter on its scheduled turn,
but remains available on every later turn until it enters.

\begin{center}
\small
\begin{tabularx}{\textwidth}{@{}p{.72in}p{.72in}Y@{}}
\toprule
Turn & Time & Arriving pieces \\
\midrule
2 & 12:30 & \textbf{Fairfax Foot:} Rigby, Dodding, Ashton, Fairfax, Bright, and Constable. \\
3 & 1:00 & \textbf{Manchester Foot:} Manchester, 1/Crfrd, and 2/Crfrd. \\
4 & 1:30 & \textbf{Leven's first group:} Reserve, Sinclair, Erskine, Dunhope, Yester, Livgstn, Coupar, and Dunfer; \textbf{all six Allied artillery pieces}. \\
5 & 2:00 & \textbf{Leven's second group:} Killhead, Cassilis, Bchlch, Loudon, Rae, Hamilton, Maitland, and Fifeshire. \\
\bottomrule
\end{tabularx}
\end{center}

\begin{scenarioBox}{Approach Column}
On the turn it enters, an ordered Foot unit has 8 MPs instead of its printed
allowance. It pays normal terrain and hexside costs, and may move through
friendly pieces normally, but it may not enter an ordered Royalist unit's ZOC.
\newline
\newline
An entering Foot unit may not attack during that turn's Allied Combat Phase.
It is otherwise an ordered unit: it exerts a ZOC and defends at full strength.
\end{scenarioBox}

\begin{scenarioBox}{Arriving Artillery}
An arriving Allied artillery piece may move up to 8 MPs during its entry March
Phase, paying Foot movement costs. It may move only through clear or hilltop
(aka rough) hexes, or along lanes, and may not enter an ordered Royalist unit's
ZOC. Once this movement ends, the piece is unlimbered and cannot move again.
\newline
\newline
An arriving artillery piece cannot bombard until a later Allied Artillery
Phase. Thus the guns entering on Turn 4 can first bombard on Turn 5, at 2 p.m.
\end{scenarioBox}

\section{Special Rules}

Except where changed below, use all standard rules from sections 1.0--15.0.

\begin{enumerate}
  \item \textbf{Royalist initiative.} Skip the Allied Rally, Artillery, March,
    and Combat Phases on Turn 1. Begin play with the Royalist Rally Phase and
    complete the rest of the turn normally. Starting with Turn 2, use the normal
    sequence of play.

  \item \textbf{Daylight.} Visibility is Clear throughout the scenario. Skip
    every Visibility Phase. The flash-rain rule and the Short Game's temporary
    cancellation of ditch effects are not used; ditch hexsides function
    normally from the beginning.

  \item \textbf{Uncommitted arrivals.} A Foot unit still off-map at the end of the
    scenario awards the Royalist player VPs equal to half its ordered combat
    strength, rounded down. Each Allied artillery piece still off-map awards 5
    VPs. These pieces do not count toward Allied demoralization.

  \item \textbf{Optional rules.} For the first playtest, do not use the Optional
    Skirmishing Rule. It can be added once the basic arrival schedule has been
    tested; it gives the exposed Allied rearguard another way to slow the
    Royalist advance and may require moving one Scottish group a turn later.
\end{enumerate}

\section{Victory}

Calculate VPs normally under 15.0, including the standard awards for the train,
artillery, leaders, demoralization, eliminated units, and disrupted Foot. Add
any VPs for Allied pieces that never entered under the special rule above.

Apply the standard victory ratios, with one scenario-specific change:

\begin{redBox}{The Royalists Must Exploit the Opportunity}
If neither player has a 1.5-to-1 ratio at the end of Turn 12, the Allied player
wins a \textbf{marginal scenario victory}. The Allied army has survived the
dangerous hours, completed its deployment, and denied Rupert the quick decision
he sought. A Royalist marginal, major, or crushing victory therefore requires
the corresponding standard ratio.
\end{redBox}

If the Allied player earns a standard victory ratio, determine its level
normally. The scenario victory described above is always marginal.


\end{document}
